\documentclass[11pt]{article}
\usepackage[left=25mm, right=25mm, top=20mm, bottom=30mm]{geometry}
\usepackage{times}
\usepackage{indentfirst}
\usepackage{titling}
\setlength{\droptitle}{-3cm}
\pretitle{\begin{center}\LARGE}
\posttitle{\par\end{center}}
\preauthor{\begin{center}\large}
\postauthor{\end{center}}
\predate{\begin{center}\normalsize}
\postdate{\end{center}\vspace{-10pt}}

\usepackage{graphicx}
\usepackage{hyperref}
\usepackage{amsmath}
\usepackage{amssymb}
\usepackage{cancel}
\usepackage{xcolor}
\usepackage{ulem}
\usepackage{soul}

\begin{document}

\section*{\normalsize{Authors' responses to the editor's comment:}}

We sincerely thank the editor for the opportunity to submit a minor revision of our manuscript.
We have carefully addressed all remaining reviewer comments while maintaining the manuscript within the four-page limit of \textit{IEEE Electron Device Letters}.
In particular, the revised manuscript now includes concise statements of the tested ranges and convergence results for the cryogenic-emulation and numerical-discretization checks, together with an explicit statement of the present validation scope.

All corresponding changes have been incorporated into the revised manuscript and are highlighted in the marked revision manuscript.
The page and line numbers reported below refer to the marked revision manuscript.

\section*{\normalsize{Authors' responses to the reviewers' comments:}}

\section*{\textit{Reviewer \# 1}}
\subsection*{\textit{Comment \# 1}}

\noindent \textit{Review wrote:}

No further comments on the manuscript.

\noindent \textit{\\Our response:}

We sincerely thank the reviewer for confirming that the previous concerns have been addressed and for having no further comments on the manuscript.

\noindent \textit{\\Corresponding change in manuscript: No}

\hspace{40pt} \noindent \textit{Location of Change: -}

\section*{\textit{Reviewer \# 2}}

We thank the reviewer for the careful reassessment of our manuscript and for recognizing that the previous concerns were addressed constructively.
We have incorporated the requested concise convergence information and explicit validation-scope statement into the main manuscript, as detailed below.

\subsection*{\textit{Comment \# 1}}

\noindent \textit{Review wrote:}

Some of the strongest supporting evidence ($\sigma_n$, temperature convergence test, discretization convergence test) appears only in the response letter.
Since some of these are key to justifying the cryogenic emulation approach and the numerical implementation, the authors should include at least a concise statement of the tested ranges and convergence results in the manuscript.

\noindent \textit{\\Our response:}

We thank the reviewer for this constructive suggestion.
We agree that the principal tested ranges and convergence outcomes should be available directly in the manuscript rather than only in the previous response letter.
Accordingly, Section~III now states that the temperature was lowered from 300~K to 83~K, the lowest convergent temperature in our TCAD setup, and that the electron capture cross section was reduced from $10^{-16}$ to $10^{-31}$~cm$^2$.
The revised text also states that the retention transient approached a limiting curve as $\sigma_n$ decreased and that its time dependence was consistent with the independently obtained temperature-lowering result.

In addition, the manuscript now reports the tested discretization ranges of 10--140 trap-energy bins and 5--50 radial CTN slices, and states that changes became negligible beyond 100 bins and 30 slices, respectively.
These concise additions place the essential convergence ranges and conclusions in the main manuscript while preserving the four-page letter format.


\noindent \textit{\\Corresponding change in manuscript: Yes \\Location of Change:}

\hspace{40pt} \noindent \textit{Section: \uppercase\expandafter{\romannumeral 3}}

\hspace{40pt} \noindent \textit{Page\# and rough location in page: Lines 4--15 and 46--48 at Page 3}

\hspace{40pt} \noindent \textit{Figure\#: -}

\subsection*{\textit{Comment \# 2}}

\noindent \textit{Review wrote:}

The scope limitation should be stated more explicitly in the main manuscript, preferably in the results section.
The present model is not validated against absolute experimental cryogenic $V_\mathrm{T}$ loss or full 4.2~K device physics.
It is important because the title and abstract may otherwise give the impression of direct cryogenic predictive capability.

\noindent \textit{\\Our response:}

We thank the reviewer for emphasizing the need to state the validation scope explicitly.
We agree that the manuscript should not imply direct validation of the absolute experimental cryogenic $\Delta V_\mathrm{th}$ magnitude or complete device physics at 4.2~K.

To prevent this interpretation, the Abstract now specifies that the TCAD simulations were calibrated \textit{only} to the time scale of the reported cryogenic retention data.
In Section~III, the TCAD results in Fig.~3 are explicitly described as a calibrated cryogenic-emulation benchmark rather than direct 4.2~K TCAD simulations.
Most importantly, the Results section now states that the present validation is limited to this calibrated cryogenic-emulation TCAD benchmark and does not establish either the absolute experimental cryogenic $\Delta V_\mathrm{th}$ magnitude or full device physics at 4.2~K.
The Conclusion was also revised to use the same calibrated cryogenic-emulation benchmark terminology.


\noindent \textit{\\Corresponding change in manuscript: Yes \\Location of Change:}

\hspace{40pt} \noindent \textit{Section: Abstract, \uppercase\expandafter{\romannumeral 3}, and \uppercase\expandafter{\romannumeral 4}}

\hspace{40pt} \noindent \textit{Page\# and rough location in page: Lines 10--13 at Page 1; Lines 24--25 and 34--40 at Page 3; Lines 76--78 at Page 3}

\hspace{40pt} \noindent \textit{Figure\#: -}

\end{document}
